% !TeX root = ../概率论与数理统计笔记.tex
% 由主文档通过 \input 载入；请编译主文档。
\chapter*{绪论：研究对象、课程总纲与学习路线}
\addcontentsline{toc}{chapter}{绪论：研究对象、课程总纲与学习路线}

\section*{从随机现象到统计推断}

\emph{概率论}（probability theory）研究随机现象的数学模型及其规律；
\emph{数理统计}（mathematical statistics）研究如何利用观测数据认识产生数据的总体，并量化结论的不确定性。
两者通过同一套随机模型相互联系：概率论通常从给定模型出发推导观测结果的分布，数理统计则从观测结果出发估计模型中的未知量、检验假设或预测新的观测。

例如，在独立重复的硬币试验中，若已知每次正面出现的概率，可以计算正面总次数的分布；
若这一概率未知，则需要根据有限次试验的结果估计它，并说明估计误差。
本书以这类从模型到数据、再从数据回到模型的问题为主线，逐步建立概念、计算方法与论证依据。

\section*{全书定位与先修知识}

本书按本科阶段的系统课程组织，主干内容覆盖概率基础、随机变量、极限定理、统计推断以及常用统计模型。
读者应具备集合与函数、一元微积分、级数和初步线性代数知识；多元积分、变量替换与矩阵二次型在使用前集中回顾。
测度论不作为起点，但会明确事件族、可测性和可积性的含义，避免把只在离散情形成立的直觉无条件推广到一般情形。

当前版本为\emph{课程总纲与分章写作大纲}：各节说明后续正文需要建立的概念、证明的结论、处理的例题与易错点。
其中少量公式用于确定记号和结论的适用范围；完整定义、证明、例题解答和习题解答将在相应章节继续展开。

\section*{课程总纲}

全书分为四篇。前三篇为主干，第四篇为选读；各章顺序反映知识依赖，并非固定课时安排。

\begin{description}
  \item[第一篇：概率模型与随机变量]
  从事件和概率出发，研究分布、数字特征、条件结构与随机变量序列的极限。
  \begin{enumerate}
    \item 第\ref{chap:ps-events}章：随机试验、事件与概率空间。
    \item 第\ref{chap:ps-conditioning}章：条件概率、独立性与贝叶斯公式。
    \item 第\ref{chap:ps-univariate}章：随机变量、分布函数与常见分布。
    \item 第\ref{chap:ps-multivariate}章：多维随机变量及其分布。
    \item 第\ref{chap:ps-moments}章：数学期望与数字特征。
    \item 第\ref{chap:ps-transformations}章：随机变量函数的分布。
    \item 第\ref{chap:ps-conditional-expectation}章：条件期望与条件方差。
    \item 第\ref{chap:ps-generating-functions}章：生成函数与特征函数。
    \item 第\ref{chap:ps-limits}章：概率不等式、大数定律与中心极限定理。
  \end{enumerate}
  \item[第二篇：数理统计的基本理论]
  由随机抽样建立统计模型，依次研究样本包含的信息、点估计、区间估计与假设检验。
  \begin{enumerate}
    \item 第\ref{chap:ps-sampling}章：总体、样本、统计量与抽样分布。
    \item 第\ref{chap:ps-information}章：充分统计量与统计信息。
    \item 第\ref{chap:ps-point-estimation}章：点估计及其评价。
    \item 第\ref{chap:ps-intervals}章：区间估计与置信区域。
    \item 第\ref{chap:ps-testing}章：假设检验的原理与方法。
    \item 第\ref{chap:ps-nonparametric}章：拟合优度、列联表与非参数方法。
  \end{enumerate}
  \item[第三篇：统计模型与数据分析]
  将估计和检验用于多个变量之间的关系与组间差异，形成建模、计算、诊断和解释的完整过程。
  \begin{enumerate}
    \item 第\ref{chap:ps-regression}章：线性回归与预测。
    \item 第\ref{chap:ps-anova}章：方差分析与试验设计初步。
  \end{enumerate}
  \item[第四篇：进一步学习的专题（选读）]
  \begin{enumerate}
    \item 第\ref{chap:ps-bayesian}章：贝叶斯统计初步。
    \item 第\ref{chap:ps-processes}章：随机过程初步。
    \item 第\ref{chap:ps-computation}章：随机模拟与统计计算。
  \end{enumerate}
\end{description}

附录\ref{app:ps-prerequisites}整理必要的数学工具，附录\ref{app:ps-reference}规划分布、推断方法与学习检查索引。
附录用于回顾与查阅，不能代替正文中对条件和推导过程的说明。

\section*{学习路径与章节依赖}

首次学习可沿第一、二篇顺序推进，再学习回归和方差分析。
第\ref{chap:ps-generating-functions}章的特征函数证明细节、
第\ref{chap:ps-information}章的完备性及信息下界、
第\ref{chap:ps-point-estimation}章的渐近理论可在第一遍阅读时暂缓，随后返回补足理论依据。
第\ref{chap:ps-limits}章会区分需要掌握的定理陈述、初等证明以及选读的证明方法。

第\ref{chap:ps-intervals}章和第\ref{chap:ps-testing}章依赖抽样分布与点估计；
第\ref{chap:ps-regression}章和第\ref{chap:ps-anova}章进一步依赖区间估计和检验。
贝叶斯统计以条件概率和似然为基础；随机过程以联合分布、条件期望与极限定理为基础；
统计计算中的独立模拟可在学完极限定理后阅读，自助法在统计推断之后阅读，马尔可夫链蒙特卡罗则在随机过程之后阅读。

每章按“问题引入---定义与条件---性质或定理---例题---习题与小结”扩写。
习题分为概念辨析、基本计算、证明与反例、综合建模四类；计算实验用于辅助理解，并明确其不能替代理论证明。

\section*{全书记号与论述约定}

\begin{description}
  \item[概率模型]
  $\Omega$ 表示样本空间，$\mathcal F$ 表示事件族，$\mathbb P$ 表示概率；事件通常记为 $A,B$。
  这些对象在第\ref{chap:ps-events}章正式定义。
  \item[随机变量与观测]
  $X,Y$ 及 $X_1,\ldots,X_n$ 表示随机变量；$x,y$ 及 $x_1,\ldots,x_n$ 表示其取值或观测值。
  $n$ 表示样本量，$\mathbf X=(X_1,\ldots,X_n)$ 表示样本向量，$\mathbf x$ 表示其观测值。
  \item[分布与数字特征]
  $F_X$、$p_X$ 和 $f_X$ 分别表示分布函数、概率质量函数和概率密度函数。
  $\mathbb E[X]$、$\operatorname{Var}(X)$ 和 $\operatorname{Cov}(X,Y)$ 分别表示期望、方差和协方差；使用前说明其存在条件。
  \item[统计模型与估计]
  $\theta$ 表示未知参数，$\Theta$ 表示参数空间，$\mathbb P_\theta$ 和 $\mathbb E_\theta$ 表示参数为 $\theta$ 时的概率与期望。
  $T(\mathbf X)$ 表示统计量；$\widehat\theta(\mathbf X)$ 表示估计量，$\widehat\theta(\mathbf x)$ 表示观测后得到的估计值。
  \item[假设与结论]
  独立性、同分布、正态性和同方差性均是需要说明的模型条件。
  精确结论、渐近结论和数值近似分别表述；不把大样本近似写成有限样本恒等式。
\end{description}


